\section{Problem Formulation}
\label{sec:problem}

This work addresses two computational problems involving least-squares
methods. The numerical data are constructed from the eight digits of
the student's enrollment number. Any zero digit is replaced by one
before constructing the matrices and vectors.

\subsection{Problem 1: QR Factorization and Linear Least Squares}

Consider the overdetermined linear system
\begin{equation}
    A\mathbf{x}\approx\mathbf{b},
\end{equation}
where $A\in\mathbb{R}^{5\times4}$, $\mathbf{x}\in\mathbb{R}^{4}$,
and $\mathbf{b}\in\mathbb{R}^{5}$. The coefficient matrix and
right-hand-side vector are defined by
\begin{equation}
    A=
    \begin{bmatrix}
        m_1+2 & m_2   & -m_3  & m_4  \\
        m_2+1 & -m_3  & m_4+2 & m_5  \\
        m_3+2 & -m_4  & m_5+1 & -m_6 \\
        -m_4  & m_5+2 & m_6   & m_7  \\
        -m_5  & m_6+1 & m_7+2 & m_8
    \end{bmatrix},
    \qquad
    \mathbf{b}=
    \begin{bmatrix}
        m_1+m_2 \\
        -m_3    \\
        m_4+m_5 \\
        -m_6    \\
        m_7+m_8
    \end{bmatrix}.
    \label{eq:system-data}
\end{equation}

The first objective is to compute the reduced QR factorization
\begin{equation}
    A=QR,
\end{equation}
where $Q\in\mathbb{R}^{5\times4}$ has orthonormal columns and
$R\in\mathbb{R}^{4\times4}$ is upper triangular.

The least-squares solution is defined as
\begin{equation}
    \mathbf{x}^{*}
    =\underset{\mathbf{x}\in\mathbb{R}^{4}}{\arg\min}
    \,\|A\mathbf{x}-\mathbf{b}\|_2^2.
    \label{eq:ls-first}
\end{equation}

The uniqueness of this solution depends on the column rank of $A$.
In particular, the solution is unique if and only if
$\operatorname{rank}(A)=4$. This condition will be assessed
computationally for the given enrollment number.

\subsection{Problem 2: Quadratic Least-Squares Fitting}

Consider six experimental data points $(x_i,y_i)$, where
\begin{equation}
    x_i=i,\qquad
    y_i=m_i\bmod 5,\qquad i=1,\ldots,6.
    \label{eq:experimental-data}
\end{equation}

The objective is to approximate these data using a quadratic
polynomial
\begin{equation}
    p(x)=ax^2+bx+c.
\end{equation}

Defining the coefficient vector
$\boldsymbol{\theta}=[a,b,c]^T$, the fitting problem can be
expressed as
\begin{equation}
    X\boldsymbol{\theta}\approx\mathbf{y},
\end{equation}
where
\begin{equation}
    X=
    \begin{bmatrix}
        1^2 & 1 & 1 \\
        2^2 & 2 & 1 \\
        3^2 & 3 & 1 \\
        4^2 & 4 & 1 \\
        5^2 & 5 & 1 \\
        6^2 & 6 & 1
    \end{bmatrix},
    \qquad
    \mathbf{y}=
    \begin{bmatrix}
        y_1 \\y_2\\y_3\\y_4\\y_5\\y_6
    \end{bmatrix}.
    \label{eq:quadratic-system}
\end{equation}

The polynomial coefficients are obtained by solving
\begin{equation}
    \boldsymbol{\theta}^{*}
    =\underset{\boldsymbol{\theta}\in\mathbb{R}^{3}}
    {\arg\min}
    \,\|X\boldsymbol{\theta}-\mathbf{y}\|_2^2.
    \label{eq:ls-quadratic}
\end{equation}

Since the six values of $x_i$ are distinct, the three columns of
$X$ are linearly independent. Therefore,
$\operatorname{rank}(X)=3$, and the least-squares solution is
unique.

Finally, the quality of the fitted polynomial is evaluated by
plotting the experimental points and the fitted curve and by
computing the mean squared error (MSE):
\begin{equation}
    \operatorname{MSE}
    =\frac{1}{6}\sum_{i=1}^{6}
    \left[y_i-p(x_i)\right]^2.
    \label{eq:mse}
\end{equation}